% Generated draft: requires mathematical and source review.
\documentclass[11pt]{article}
\usepackage[margin=25mm]{geometry}
\usepackage{fontspec}
\setmainfont{Latin Modern Roman}
\usepackage{amsmath,amssymb,mathtools,unicode-math}
\setmathfont{Latin Modern Math}
\usepackage{enumitem,needspace,xurl,hyperref}
\setlength{\parindent}{0pt}
\setlength{\parskip}{5pt}
\setlength{\emergencystretch}{3em}
\newcommand{\N}{\mathbb{N}}
\newcommand{\Z}{\mathbb{Z}}
\newcommand{\Q}{\mathbb{Q}}
\newcommand{\R}{\mathbb{R}}
\newcommand{\C}{\mathbb{C}}
\newcommand{\F}{\mathbb{F}}
\newcommand{\A}{\mathbb{A}}
\newcommand{\PP}{\mathbb{P}}
\newcommand{\E}{\mathbb{E}}
\newcommand{\eps}{\varepsilon}
\newcommand{\dd}{\,\mathrm{d}}
\begin{document}
\section*{UnsolvedMath ID: 30006085}
\textbf{Source problem code:} catalogue entry
\subsection{Definitions and mathematical statement}
The supplied record asks for a connection between an invariant measure of the renormalized stochastic geometric heat equation and Brownian loop measure. Its displayed regularized equation can be transcribed schematically as
\[
\begin{aligned}
\partial_t u_\varepsilon^\alpha
={}&\partial_x^2u_\varepsilon^\alpha
+\Gamma^\alpha_{\beta\gamma}(u_\varepsilon)
\partial_xu_\varepsilon^\beta\,\partial_xu_\varepsilon^\gamma
+h^\alpha(u_\varepsilon)
+\sum_i\sigma_i^\alpha(u_\varepsilon)\xi_\varepsilon^i\\
&+\sum_{\tau\in B_\xi}C_\varepsilon(\tau)
\Upsilon^\alpha_{\Gamma,\sigma}[\tau](u_\varepsilon).
\end{aligned}
\]
Here \(\alpha,\beta,\gamma\) are coordinate indices, with repeated \(\beta,\gamma\) summed; \(\Gamma\), \(h\), and \(\sigma_i\) denote the geometric coefficients, drift, and noise coefficients appearing in the source. The symbols \(\xi_\varepsilon^i\), \(C_\varepsilon(\tau)\), \(B_\xi\), and \(\Upsilon_{\Gamma,\sigma}[\tau]\) denote the regularized noises, renormalization constants, counterterm index set, and counterterm expressions, respectively. Their precise definitions and admissibility conditions are not supplied. This transcription repairs the corrupted placement of indices and summation signs; it does not specify a complete stochastic equation.

For a Markov process on a specified measurable state space \(\mathcal X\), write \(P_tF(v)\) for the expectation of \(F\) at time \(t\), conditional on initial state \(v\). A probability measure \(\mu\) is invariant if
\[
\int_{\mathcal X}P_tF(v)\,\mu(\dd v)
=\int_{\mathcal X}F(v)\,\mu(\dd v)
\quad\text{for every }t\geq0
\text{ and every bounded measurable }F:\mathcal X\to\R.
\]
If the intended Brownian loop measure is instead an infinite measure, invariance must be formulated for nonnegative measurable functions, allowing extended integrals, and any normalization must be specified separately.

\textbf{Requested problem, with unresolved qualifications.} Identify and prove the expected relationship between an invariant measure of the limiting renormalized geometric heat flow and an appropriate Brownian loop measure, in the setting described by the supplied record as the ``full subcritical regime.'' If the intended relationship is equality with a Brownian loop probability measure \(\mu_{\mathrm{BL}}\), the target would be the displayed invariance identity with \(\mu=\mu_{\mathrm{BL}}\). This is only a conditional formulation: the supplied excerpt says that the connection remains open, but does not specify equality, normalization, or any alternative relationship.

The target manifold, spatial domain, boundary conditions, noise law and regularity, drift assumptions, renormalization prescription, limiting Markov process, and convention for Brownian loop measure are absent. Consequently, the record does not determine a precise universally quantified conjecture. In particular, the excerpt's assertion of geometric renormalization in the full subcritical regime does not itself establish that Brownian loop invariance is expected throughout that regime.
\subsection{Short English statement}
The problem asks how Brownian loop measure is related to an invariant measure of a renormalized stochastic geometric heat flow. The supplied record labels this connection open. A precise statement requires the geometric setting, noise, drift, renormalization, boundary conditions, and loop-measure convention; the excerpt does not provide these or specify the exact relationship to prove.
\subsection{Sources}
\begin{itemize}
\item Supplied local UnsolvedMath registry record, permanent ID \(30006085\), source problem code \texttt{catalogue entry}, title ``Brownian Loop Measure for Stochastic Geometric Heat Flow.'' No individual proposer is supplied; the proposed year is recorded as 2024.
\item Supplied source citation: \emph{Directions in Rough Analysis} (2025), Oberwolfach Reports, pp.\ 2901--2948, DOI \texttt{10.4171/owr/2024/49}. \url{https://doi.org/10.4171/owr/2024/49}. The literature-source list labels this ``Oberwolfach Report 49/2024.'' The excerpt attributes its geometric renormalization theorem to reference [4], the space-time-white-noise result to [5], and the Regularity Structures solution theory to [11, 7, 9, 2]; bibliographic entries for these numbered references are not supplied.
\item Supplied literature source: ``ANR project on low-regularity dynamics.'' \url{https://anr.fr/Project-ANR-23-ERCS-0011}.
\item Literature assessment dated exactly \texttt{2026-08-22}: the supplied curation labels the problem open, reports that Brownian loop measure is expected to be invariant, says the primary source identifies the connection as open, and reports that no later proof was found. This is supplied dated curation, not an independently verified literature assessment. Neither linked source was visited for this draft.
\item The registry labels the statement \texttt{corrected\_verified} and describes its changes as repairs to delimiters, notation, grammar, and typography. These supplied labels do not establish mathematical correctness or resolve the missing hypotheses.
\end{itemize}
\end{document}
